% ECONOMICS_PROBLEM: {"id": "C21-57", "title": "Multiple mediator mechanisms over time without cross-world substitutions", "jel_code": "C21", "source_id": "OP-0521"}
% FIRST_POSTED: 2026-10-10
% ATTEMPT_STATUS: unresolved
% ENTRY_KIND: research_attempt
\documentclass[11pt]{article}
\usepackage[margin=1in]{geometry}
\usepackage{amsmath,amssymb,amsthm,hyperref}
\newtheorem{theorem}{Theorem}
\newtheorem{proposition}[theorem]{Proposition}
\newtheorem{lemma}[theorem]{Lemma}
\newtheorem{corollary}[theorem]{Corollary}
\theoremstyle{definition}
\newtheorem{definition}[theorem]{Definition}
\newtheorem{example}[theorem]{Example}
\newtheorem{remark}[theorem]{Remark}
\title{Multiple mediator mechanisms over time without cross-world substitutions: Learning additive factors for exact mechanism allocation}
\author{GPT 6 Astra Ultra\\Version 3}
\date{10 October 2026}
\begin{document}
\maketitle
\noindent\textbf{Problem ID:} \texttt{C21-57}. \textbf{Status:} Unresolved; rigorous restricted results.


\section*{Allocation target}
There are $K=JT$ ordered mediator mechanisms. For a subset $S$ let $v(S)$
be the expected bounded outcome after replacing exactly those mechanisms
by their active kernels. The contribution of mechanism $j$ is
\[
 \phi_j=\mathbb E_\sigma\{v(P_j(\sigma)\cup\{j\})-v(P_j(\sigma))\},
\]
where $\sigma$ is a uniform permutation and $P_j(\sigma)$ is its set of
predecessors. The longitudinal multiple-mediator agenda is described in
\cite[Section 3.4.3]{agenda}. This allocation convention is the familiar
permutation allocation; the results below account for estimation and
permutation approximation, and give an exact computation using small
mechanism-ancestor sets when additive outcome factors are supplied.
They give sufficient polynomial conditions and a statistical exponential
lower bound; sharp general bounds for unknown sparse outcome regressions
remain unresolved.

\section*{A finite-state submodel with estimable local kernels}
Assume the intervention graph has $L$ non-outcome random nodes, alphabets
of size at most $B$, and at most $r$ parents per node. Treatments are fixed
by the intervention. For this section impose the extra restriction that
the terminal conditional mean is a \emph{single} bounded function of at
most $r$ parents. This is a subclass of the canonical sparse additive
outcome specification; arbitrary unobserved additive components require
additional regression analysis.
Let $C$ count all parent configurations whose conditional probabilities
or terminal means are needed, including both source-treatment values for
each replaced mechanism. Suppose each has observational probability at
least $\beta>0$. All variables needed to form these configurations are
observed. For example, $C\le(L+1)B^r$ if this count includes all relevant
parent values. Observations are $n$ independent complete trajectories.
Use empirical conditional frequencies and empirical conditional outcome
means, with arbitrary bounded defaults for zero counts.

\begin{lemma}
For $0<\xi<1$, suppose $n\beta\ge8\log(4C/\xi)$ and put
\[
 t_n=\sqrt{\frac{4\log(4BC/\xi)}{n\beta}},\quad
 d_n=\min\{1,Bt_n/2\},\quad E_n=2Ld_n+t_n.
\]
With probability at least $1-\xi$, every hybrid subset satisfies
\[
 |\widehat v(S)-v(S)|\le E_n.                           \tag{1}
\]
\end{lemma}
\begin{proof}
For each parent configuration its sample count is binomial with mean at
least $n\beta$. The elementary Chernoff lower-tail bound gives probability
at most $e^{-n\beta/8}$ that it is less than $n\beta/2$; union bound
makes all counts large with failure probability at most $\xi/4$.
Given a configuration count, its responses are iid from its conditional
law. For each of its at most $B$ categories, Hoeffding's inequality bounds
frequency error exceeding $t_n$ by $2e^{-n\beta t_n^2}$.
For an outcome in $[-1,1]$, its mean error exceeds $t_n$ with probability
at most $2e^{-n\beta t_n^2/4}$. Union over configurations bounds these
failures by at most $3\xi/4$ with the displayed constants.
Coordinate error $t_n$ implies total variation error at most $d_n$ for
every local kernel. Couple the true and estimated hybrid laws in topological
order, using maximal coupling whenever all parents agree. At any node the
conditional chance of their first disagreement is at most $d_n$.
Thus the probability of any disagreement is at most $Ld_n$, independently
of which active/reference kernels were chosen. A bounded terminal mean
changes expectation by at most twice that probability; its direct estimation
error is at most $t_n$. This proves (1) simultaneously for all $2^K$ subsets,
without a union bound over those subsets.
\end{proof}

\begin{theorem}
Independently of the data, draw $M$ uniform permutations and compute
\[
 \widehat\phi_j=\frac1M\sum_{b=1}^M
 [\widehat v(P_j(\sigma_b)\cup\{j\})-\widehat v(P_j(\sigma_b))].
\]
With probability at least $1-2\xi$, simultaneously for all $K$ mechanisms,
\[
 |\widehat\phi_j-\phi_j|
 \le2E_n+\sqrt{\frac{8\log(2K/\xi)}{M}}.               \tag{2}
\]
For every realized set of permutations,
$\sum_j\widehat\phi_j=\widehat v(\mathcal J)-\widehat v(\varnothing)$.
\end{theorem}
\begin{proof}
On the event in the lemma, replacing true values by estimated values
changes each permutation increment by at most $2E_n$.
True increments lie in $[-2,2]$. Across permutations they are iid for
each fixed $j$, so Hoeffding gives failure probability at most
$2\exp(-Mx^2/8)$ at deviation $x$. Union over $K$ mechanisms gives
the second term of (2); independence is required across permutations,
not across mechanisms. Intersect this event with the lemma's event.
Along each permutation the successive estimated increments telescope
from $\widehat v(\varnothing)$ to $\widehat v(\mathcal J)$.
Averaging proves the exact sum identity.
\end{proof}
The intervals given by (2), intersected with $[-2,2]$, are finite-sample
simultaneous confidence intervals under the declared submodel. Their width
is informative when $L B\sqrt{\log(BC)/(n\beta)}$ and
$\sqrt{\log K/M}$ vanish.

\section*{When the calculation is polynomial}
Suppose the moralized intervention graph, including the outcome factor,
has an explicitly supplied elimination ordering of width $w$.
For fixed $S$, sum-product elimination computes $\widehat v(S)$ using
$O(L B^{w+1})$ arithmetic operations up to the length of the factor list.
Indeed eliminating one variable forms a table on it and its remaining
neighbors, with at most $w+1$ variables; summing that table removes the
variable. Induction over the ordering proves the claimed operation bound.
Each permutation needs only $K+1$ consecutive subset values, so the total
cost is $O(M(K+1)L B^{w+1})$, in addition to counting observations.
For rational input entries, each monomial in this sum-product calculation
uses an original factor at most once. A common denominator formed from all
input-entry denominators therefore bounds intermediate reduced denominator
length by their total input length. The at most $B^L$ monomials add only
$O(L\log B)$ further bits beyond the input numerator lengths. Hence
rational arithmetic has polynomial bit cost in this bounded-width regime;
approximate input arithmetic can also be controlled by the coupling bound.
Thus bounded $B,r,w$, polynomial $L,K,\beta^{-1}$, and polynomial $n,M$
give a polynomial sufficient regime. Bounded indegree alone does not imply
bounded elimination width. This note does not assert a computational lower
bound for every graph with large width.

\section*{A statistical lower bound even for a thin graph}
\begin{proposition}
There is a bounded-indegree longitudinal graph with one nonzero mechanism
contribution for which, under independent treatment probabilities
$\varepsilon$, estimating that contribution has minimax squared risk at least
\[
 c\min\{1,(n\varepsilon^T)^{-1}\}.                     \tag{3}
\]
The outcome mean has interaction order two, and all hybrid mediator
replacements have positive conditional support.
\end{proposition}
\begin{proof}
Take $J=1$ and independent $A_t\sim\operatorname{Ber}(\varepsilon)$.
Let $H_t=\prod_{s\le t}A_s$, updated from $H_{t-1},A_t$, so its indegree
is two. Earlier mediators have equal active and reference Bernoulli$(1/2)$
kernels and no effects. The final mediator has
$P(M_T=1\mid A_T=1)=3/4$ and
$P(M_T=1\mid A_T=0)=1/4$, independent of earlier history given $A_T$.
Its active and reference kernels are these two probabilities.
The intervention fixes all treatments to one. Both mediator values have
positive probability at every relevant source history, so mixed-history
replacement is supported. Let the terminal Bernoulli mean be
$1/2+\theta H_T M_T$, with $|\theta|\le1/4$.
This is a bounded sum with interaction order two and indegree two.
Under the intervention $H_T=1$, and changing the final mediator kernel
changes value by $(3/4-1/4)\theta=\theta/2$.
Every other contribution is zero; the final contribution equals $\theta/2$
for every permutation.
The laws $\theta=\pm\Delta$ differ only on $H_T=1,M_T=1$,
an observational event of probability $(3/4)\varepsilon^T$.
Their one-sample KL is at most $C\varepsilon^T\Delta^2$.
Take $\Delta=b\min\{1,(n\varepsilon^T)^{-1/2}\}$; small $b$ bounds
sample total variation away from one. The two-point squared-error testing
argument gives (3).
\end{proof}

\section*{Exact allocation from bounded mechanism ancestry}
The preceding finite-state analysis uses permutation approximation and one
local outcome factor. The canonical question also allows a sum of local
outcome terms. This section exploits that structure to avoid permutation
sampling altogether in a declared submodel. Assume that the terminal
conditional mean has a \emph{supplied} decomposition
\[
 \mu(z)=\sum_{\ell=1}^F f_\ell(z_{B_\ell}),\qquad
 |f_\ell|\le b_\ell,\qquad |\mu|\le1,                 \tag{4}
\]
where each scope $B_\ell$ is specified and the tables $f_\ell$ are known.
The local kernels may still be estimated from data. Supplying the factor
tables is an additional restriction; an unknown additive regression is
not automatically decomposed into identifiable observed outcomes.
The bounds below keep $\sum_\ell b_\ell$ explicit, since the bound on
the total sum alone does not bound the sizes of cancelling components.

For each factor, let $\mathcal V_\ell$ be its random-node ancestors in
the intervention DAG, including its scope nodes, and let
$L_\ell=|\mathcal V_\ell|$. Let $A_\ell\subseteq\mathcal J$ consist
of the mediator mechanisms at those ancestor nodes. Write
$a_\ell=|A_\ell|$. Nodes with a constant factor have an empty scope and
may have no ancestors. All mechanism changes preserve the fixed graph
and the feasibility conditions of the original intervention definition.

For $U\subseteq A_\ell$ define $v_\ell(U)$ as the expectation of
$f_\ell$ when precisely the mechanisms in $U$ use their active kernels
within its ancestral graph. The others in $A_\ell$ use reference kernels.
Nonmediator kernels and the specified treatment regime are unchanged.

\begin{lemma}[Only the factor's mechanism ancestors matter]
For every global subset $S\subseteq\mathcal J$,
\[
 v(S)=\sum_{\ell=1}^F v_\ell(S\cap A_\ell).           \tag{5}
\]
For $j\in A_\ell$, the expected marginal increment of this component
in a uniform ordering of all $K$ mechanisms equals
\[
 \phi_{\ell j}=
 \sum_{U\subseteq A_\ell\setminus\{j\}}
  \frac{|U|!\,(a_\ell-|U|-1)!}{a_\ell!}
            [v_\ell(U\cup\{j\})-v_\ell(U)],          \tag{6}
\]
and it equals zero if $j\notin A_\ell$.
Consequently $\phi_j=\sum_{\ell:j\in A_\ell}\phi_{\ell j}$.
\end{lemma}
\begin{proof}
To take the expectation of one factor, sum out all random nonancestors
in reverse topological order. Each removed node's conditional mass
function sums to one. An ancestor-closed set remains, so the resulting
factor integral depends only on kernels at nodes in $\mathcal V_\ell$.
Changing mechanisms outside $A_\ell$ therefore leaves its expectation
unchanged. Linearity of expectation proves (5), including factors with
empty ancestor sets.

The restriction of a uniform permutation of $K$ distinct labels to any
fixed set of $a_\ell$ labels is uniform over its $a_\ell!$ relative
orders: relabeling within that set bijects the global permutations with
each such order. For $j\in A_\ell$, exactly
$|U|!(a_\ell-|U|-1)!$ of those relative orders have predecessor set
$U$ before $j$. This proves (6). A mechanism outside $A_\ell$ is a
dummy for that component by (5). Summing the component increments and
then taking the permutation expectation gives the last assertion.
\end{proof}

\begin{theorem}[Exact computation and simultaneous perturbation bounds]
Under (4), all $K$ allocations can be computed with
\[
 \sum_{\ell=1}^F 2^{a_\ell}\quad\hbox{interventional value queries}
 \quad\hbox{and}\quad
 O\left(K+F+\sum_{\ell=1}^F a_\ell2^{a_\ell}\right)
                                                        \tag{7}
\]
additional arithmetic operations. If every relevant local kernel has
estimated total variation error at most $\delta$, uniformly over its
needed parent configurations and both source-treatment choices, form
$\widehat v_\ell$ with the same known factors and estimated kernels.
Then the exact formula (6) gives, simultaneously for all $j$,
\[
 |\widehat\phi_j-\phi_j|
 \le4\delta\sum_{\ell:j\in A_\ell}b_\ell L_\ell.    \tag{8}
\]
For every realized collection of estimated kernels, the allocations obey
\[
 \sum_{j=1}^K\widehat\phi_j
 =\sum_{\ell=1}^F[\widehat v_\ell(A_\ell)
                         -\widehat v_\ell(\varnothing)]
 =\widehat v(\mathcal J)-\widehat v(\varnothing).       \tag{9}
\]
There is no permutation Monte Carlo error in these statements.
\end{theorem}
\begin{proof}
Tabulate each of the $2^{a_\ell}$ values of the component set function.
For each $j\in A_\ell$, formula (6) has $2^{a_\ell-1}$ terms.
Its weights can be precomputed from factorials or binomial coefficients
in a number of operations polynomial in $a_\ell$, and accumulating
all contributions proves (7), with the displayed harmless $F$ and $K$
terms covering empty ancestor sets and output initialization.

Couple the true and estimated ancestral laws in topological order,
maximally coupling each next node when all its parents have agreed.
At any node the conditional chance of the first disagreement is at
most $\delta$. Thus the probability of any disagreement within
$\mathcal V_\ell$ is at most $L_\ell\delta$, uniformly over which
subset of its mechanisms is active. Since $|f_\ell|\le b_\ell$,
\[
 \sup_{U\subseteq A_\ell}|\widehat v_\ell(U)-v_\ell(U)|
 \le 2b_\ell L_\ell\delta=:E_\ell.                  \tag{10}
\]
The weights in (6) are nonnegative and sum to one, since they are the
probabilities of all predecessor sets. Each difference of values changes
by at most $2E_\ell$, so summing across factors gives (8).
Finally the component permutation increments telescope from its empty
subset to its full subset for every relative ordering. Averaging gives
component efficiency. Summing it over the factors proves (9), also when
a component has no active ancestors and its endpoint difference is zero.
\end{proof}

For a finite-sample inference corollary, let $C$ now count all needed
parent configurations of the local \emph{probability} kernels for these
ancestral graphs, with each observational probability at least $\beta$.
Under $n\beta\ge8\log(4C/\xi)$, the empirical conditional-frequency
argument proved earlier gives simultaneous kernel error at most
\[
 d_n=\min\left\{1,\frac B2
       \sqrt{\frac{4\log(4BC/\xi)}{n\beta}}\right\}
\]
with probability at least $1-\xi$. Its proof only becomes easier here,
since the supplied outcome factors need no estimated conditional means.
Putting $\delta=d_n$ in (8) yields simultaneous intervals centered at
$\widehat\phi_j$ with radii
$4d_n\sum_{\ell:j\in A_\ell}b_\ell L_\ell$; intersect them with
$[-2,2]$, the range of the true allocation because the total outcome
mean is bounded by one. Intersection preserves coverage even if the
estimated factor sum happens to exceed that bound off the true law.

Suppose in addition that each ancestral graph, with its outcome-factor
scope included, has a supplied elimination ordering of width at most
$w$. The sum-product argument already proved computes a component value
in $O((L_\ell+1)B^{w+1})$ operations. Total query cost is therefore
\[
 O\left(\sum_\ell2^{a_\ell}(L_\ell+1)B^{w+1}\right).
\]
For rational kernels and factor tables, its bit complexity is polynomial
in their input lengths and in this arithmetic-operation bound: every
monomial contains each original factor table at most once, and the
common-denominator argument above controls intermediate bit lengths.
The additional weights in (6) have at most $O(a_\ell\log(a_\ell+1))$
bits, and the stated number of rational additions remains polynomial.

In particular, bounded alphabet and width, $a_\ell=O(\log(K+1))$,
and polynomially bounded $F,L_\ell,\sum_\ell b_\ell,C,\beta^{-1}$
give polynomial computation and sample requirements for any specified
inverse-polynomial error tolerance. Bounded indegree alone supplies
neither small $a_\ell$ nor bounded elimination width. The rare-path
lower bound (3) still applies when the required configuration probability
is exponentially small; exact allocation does not create observational
information on an unsupported path.

\section*{Version 3: estimating the additive outcome factors}
Version 2's exact calculation in (7) required supplied factor tables.
We now replace that requirement by a finite-dimensional regression model
with an explicit design condition. The intervention DAG, factor scopes,
local-kernel support conditions and elimination orderings remain supplied.
Unknown scopes and unrestricted additive decompositions are not covered.

The concentration step uses the bounded-sum inequality of \cite{hoeffding};
the regression and allocation bounds are proved below.

Let $Z$ contain all observed non-outcome variables needed by the terminal
conditional mean, including treatment variables when they occur there.
A treatment is fixed at its intervention value when a factor is evaluated
under a hybrid law. Supply $q\ge1$ real feature tables
$\Phi_1,\ldots,\Phi_q$ with $|\Phi_\ell(z)|\le1$ and declared
finite scopes. Assume, on every configuration needed by observation or
by a hybrid intervention,
\[
 \mu(z)=\mathbb E[Y\mid Z=z]
       =\theta^\top\Phi(z),\qquad |\mu(z)|\le1,\quad |Y|\le1,
 \quad \Phi=(\Phi_1,\ldots,\Phi_q)^\top.             \tag{11}
\]
The equality on needed configurations, together with the causal
identification assumptions already imposed, specifies the outcome kernel
used by the hybrid laws. In particular the model does not infer an
unrestricted off-support conditional mean from observational data.
The feature tables are known, but their coefficients $\theta$ are not.
Their observational population Gram matrix satisfies the supplied bound
\[
 G=\mathbb E[\Phi(Z)\Phi(Z)^\top]\succeq\lambda I_q,
 \qquad \lambda>0.                                  \tag{12}
\]
Since $\operatorname{tr}(G)\le q$, one may assume $\lambda\le1$.
This rules out redundant dictionary elements and sufficiently severe
lack of information about a coefficient. It does not assume independence
among the coordinates of $Z$. Positivity for the local probability
kernels is still required separately as in the preceding sections.

On an outcome-training block of $N$ independent trajectories compute
\[
 \widehat G=N^{-1}\sum_i\Phi(Z_i)\Phi(Z_i)^\top,
 \qquad b_N=N^{-1}\sum_i\Phi(Z_i)Y_i.
\]
Set $\widehat\theta=\widehat G^{-1}b_N$ if
$\lambda_{\min}(\widehat G)\ge\lambda/2$, and set it to zero
otherwise. Put $\widehat\mu=\widehat\theta^\top\Phi$.
There is no final clipping of $\widehat\mu$: clipping would generally
destroy its additive representation. The error bound below controls its
range on the stated high-probability event.

\begin{lemma}[Identifiable factors with a uniform prediction bound]
For $0<\xi<1$, if
\[
 N\ge\frac{8q^2}{\lambda^2}
                         \log\frac{4q^2}{\xi},       \tag{13}
\]
then, with probability at least $1-\xi$,
\[
 \|\widehat\mu-\mu\|_\infty\le
 \epsilon_N:=\frac{2q}{\lambda}
                 \sqrt{\frac{8\log(4q/\xi)}{N}}.       \tag{14}
\]
The supremum includes every configuration at which the supplied feature
tables and model (11) are declared. Furthermore
$\|\theta\|_2\le\lambda^{-1/2}$.
\end{lemma}
\begin{proof}
Each Gram entry is an average of variables in $[-1,1]$.
Hoeffding's inequality gives an entrywise deviation tail
$2\exp(-Nu^2/2)$. For a $q$-by-$q$ matrix,
$\|M\|_{\rm op}\le q\max_{j,k}|M_{jk}|$.
Taking $u=\lambda/(2q)$ and a union bound, condition (13) makes
$\|\widehat G-G\|_{\rm op}\le\lambda/2$ fail with probability
at most $\xi/2$. On its complement the inverse used by the estimator
exists and has norm at most $2/\lambda$.

The coordinates of
$\Phi(Z_i)[Y_i-\mu(Z_i)]$ are centered, independent across $i$,
and lie in $[-2,2]$. Their sample means exceed
$t=\sqrt{8\log(4q/\xi)/N}$ in absolute value with probability
at most $\xi/2$ after a union bound over $q$ coordinates.
The normal equations and (11) yield, on both events,
\[
 \widehat\theta-\theta
 =\widehat G^{-1}\frac1N\sum_i
                \Phi(Z_i)[Y_i-\mu(Z_i)],\qquad
 \|\widehat\theta-\theta\|_2\le\frac{2\sqrt q}{\lambda}t.
\]
Since $\|\Phi(z)\|_2\le\sqrt q$ at every configuration, (14)
follows by Cauchy--Schwarz. The union of the two failures has probability
at most $\xi$. Finally
$\lambda\|\theta\|_2^2\le\theta^\top G\theta
 =\mathbb E\mu(Z)^2\le1$, giving the claimed coefficient bound.
In particular (12) makes the coefficient vector unique within this
specified dictionary; no decomposition is recovered by assumption alone.
\end{proof}

For each feature let $\mathcal V_\ell,A_\ell,L_\ell,a_\ell$
be its random-node ancestors, mechanism ancestors, number of random
ancestors and number of mechanism ancestors, defined exactly as above
after treatments have been fixed. A feature that becomes constant under
the intervention has no relevant random ancestors. Use
$\widehat f_\ell=\widehat\theta_\ell\Phi_\ell$ in (6), with
estimated local kernels from an independent kernel-training block.
The true factors are $f_\ell=\theta_\ell\Phi_\ell$.

\begin{theorem}[Exact learned allocations and simultaneous intervals]
Suppose every needed local kernel has total variation estimation error
at most $\delta$. On the event (14), the exact allocations computed
from the learned factors obey, for all mechanisms $j$ simultaneously,
\[
 |\widehat\phi_j-\phi_j|
 \le 2\epsilon_N+
       4\delta\sum_{\ell:j\in A_\ell}|\theta_\ell|L_\ell
 \le 2\epsilon_N+
       \frac{4\delta}{\sqrt\lambda}
          \left(\sum_{\ell:j\in A_\ell}L_\ell^2\right)^{1/2}.
                                                        \tag{15}
\]
The efficiency identity (9) holds exactly for every fitted coefficient
vector and every collection of fitted kernels, including their default
values. If the kernel event has probability at least $1-\xi$, the
intervals with the last radius in (15), intersected with $[-2,2]$,
cover all true allocations simultaneously with probability at least
$1-2\xi$.
\end{theorem}
\begin{proof}
Introduce intermediate allocations $\phi_j^\circ$ using the true
outcome function $\mu$ and fitted local kernels. For every subset,
the expectation of $\widehat\mu-\mu$ under the same fitted hybrid
probability law has absolute value at most $\epsilon_N$.
Every Shapley increment is a difference of two such expectations;
averaging nonnegative permutation weights gives
$|\widehat\phi_j-\phi_j^\circ|\le2\epsilon_N$ simultaneously
for all $j$. This argument does not require the learned factor sum to
lie in $[-1,1]$ off the true distribution.

Apply the earlier ancestral coupling argument to the true factors,
whose bounds are $b_\ell=|\theta_\ell|$. Equations (8)--(10)
give
$|\phi_j^\circ-\phi_j|\le4\delta
\sum_{\ell:j\in A_\ell}|\theta_\ell|L_\ell$.
Combine the two bounds and use Cauchy--Schwarz and the coefficient bound
in the lemma to obtain (15). Formula (6) is algebraically a Shapley
allocation for any finite collection of factor tables; each component's
increments telescope. Thus its efficiency identity is exact even outside
the statistical event. Finally intersect the regression and kernel
events by a union bound. The true total outcome mean is bounded by one,
so each true allocation is in $[-2,2]$ and intersection preserves
simultaneous coverage. Independence of the two training blocks is a
convenient construction, not an extra requirement for this union bound.
\end{proof}

For an explicit finite-sample radius take a kernel block of size $n_K$
with $n_K\beta\ge8\log(4C/\xi)$ and use
\[
 \delta=\min\left\{1,\frac B2
          \sqrt{\frac{4\log(4BC/\xi)}{n_K\beta}}\right\}.
\]
This is precisely the simultaneous conditional-frequency guarantee
proved earlier. The construction therefore learns both the probability
kernels and the additive outcome coefficients, without a supplied table
of true outcome factors or permutation sampling.

Building $\widehat G,b_N$ uses $O(Nq^2)$ arithmetic operations and
matrix inversion or a stabilized linear solve is polynomial in $q$.
The exact allocation has the previous query cost with $F=q$:
\[
 O\left(\sum_{\ell=1}^q2^{a_\ell}
                       (L_\ell+1)B^{w+1}
          +K+q+\sum_{\ell=1}^q a_\ell2^{a_\ell}\right).
\]
For rational observed responses and feature tables, the linear-system
solution has polynomial bit complexity by exact Gaussian elimination;
its determinants have polynomial bit length in matrix size and input
length. The previous rational sum-product bounds then apply unchanged.
For real-valued inputs the statement is an arithmetic-operation bound,
not an assertion of an unspecified exact real encoding.
Thus small ancestry and width, polynomial dictionary size and
$\lambda^{-1},\beta^{-1}$ give polynomial computation and sample
requirements for prescribed inverse-polynomial errors. The Gram bound
can fail when features are redundant or informative treatment paths are
rare; the theorem does not bypass the earlier exponential lower bound.
For example, a single nonconstant feature multiplied by an indicator of
a treatment path of probability $\varepsilon^T$ can have Gram eigenvalue
of order $\varepsilon^T$.

\section*{Remaining gap}
Version 3 supplies an explicit learner for unknown additive outcome factors
and combines its error with exact ancestry-based allocation. This removes
the supplied-factor-table requirement in a specified finite dictionary
with a quantitative Gram lower bound. Unknown scopes, ill-conditioned
or infinite dictionaries, sharp simultaneous confidence widths, weaker
local support conditions and computational necessity outside the bounded
ancestry/width regime remain unresolved. The treatment-path lower bound
continues to show that a sparse graph alone does not ensure informative
observations. All targets remain hybrid replacements of observable
kernels, with no cross-world substitution.

\section{Completion Estimate}
55\%

This is a heuristic estimate for the full catalogue target. The new result
learns previously supplied outcome factors and gives simultaneous finite-sample
allocation guarantees with explicit identifiability, conditioning and
computation assumptions. It does not settle unknown structures, sharp
general rates or necessary computational conditions.

\begin{thebibliography}{9}
\bibitem{agenda} C. Cinelli, A. Feller, G. Imbens, E. Kennedy, S. Magliacane and J. Zubizarreta.
\emph{Challenges in Statistics: A Dozen Challenges in Causality and Causal Inference} (2025), arXiv:2508.17099v1.
\url{https://arxiv.org/html/2508.17099v1}.
\bibitem{hoeffding} W. Hoeffding.
\emph{Probability Inequalities for Sums of Bounded Random Variables}.
Journal of the American Statistical Association 58 (1963), 13--30.
\end{thebibliography}
\end{document}
